\chapter{Eosin-5-Maleimide (EMA) Binding Test}

\section*{What Is This Test?}
The eosin-5-maleimide binding test measures fluorescent dye binding to red-cell membrane proteins. Reduced binding is a useful laboratory marker for hereditary spherocytosis and some related red-cell membrane disorders.

\section*{Alternative Names}
EMA Binding Test; Eosin-5-Maleimide Test; Hereditary Spherocytosis Screen.

\section*{Why This Test Is Ordered}
Testing is used when chronic hemolysis, anemia, jaundice, splenomegaly, increased mean corpuscular hemoglobin concentration, or spherocytes on blood film suggests hereditary spherocytosis. It is especially useful when the diagnosis is uncertain after the CBC, reticulocyte count, smear, bilirubin, LDH, haptoglobin, and direct antiglobulin test.

\section*{How This Test Is Performed}
Patient red blood cells are incubated with fluorescent eosin-5-maleimide, washed, and analyzed by flow cytometry. Fluorescence is compared with reference cells and appropriate controls.

\section*{Preparation, Risks, and Limitations}
No fasting is required. Risks are limited to venipuncture. Recent transfusion, severe reticulocytosis, autoimmune hemolysis, and some other membrane disorders can complicate interpretation. A normal result does not exclude every hereditary hemolytic anemia.

\section*{Understanding Results}
Reduced EMA fluorescence supports hereditary spherocytosis in the correct clinical setting. Results are interpreted with evidence of hemolysis and family history; genetic testing or specialized membrane studies may be needed in atypical cases.

\section*{References}
International Council for Standardization in Haematology guidance on laboratory diagnosis of hereditary spherocytosis.

\clearpage
\section*{Understanding Results and Follow-up}
The laboratory report should identify the specimen, method, units, reference interval or decision limit, and whether the result is positive, negative, abnormal, or inconclusive. Reference intervals vary with age, sex, pregnancy, specimen type, assay, and laboratory; a result outside a range is not automatically a diagnosis, and a result within a range does not always exclude disease. Discuss unexpected results with the ordering clinician, who can decide whether repeat or confirmatory testing, treatment, or referral is appropriate. Do not change medicines or delay urgent care based on an isolated result.


\section*{Regulatory Status and Authorization Date}
Regulatory status applies to the specific assay, instrument, manufacturer, and intended use rather than to the test name alone. No single approval date can be assigned to this test category. For a commercial assay, verify the current product in the relevant regulator's database: the U.S. Food and Drug Administration (FDA) clearance, approval, or authorization; India's Central Drugs Standard Control Organisation (CDSCO) licence or permission; the European Union In Vitro Diagnostic Regulation (EU IVDR) CE certificate issued through the applicable conformity-assessment route; or the United Kingdom Medicines and Healthcare products Regulatory Agency (MHRA) registration and UKCA/CE status. Laboratory-developed tests may not have an individual FDA, CDSCO, EU IVDR, or MHRA approval date.
\clearpage
